\section{不定積分}
以下は,各被積分関数が定義される1つの区間上で用いる.\,$C$はその区間での任意定数である.\,$a>0$,\,$b\ne0$とする.
\begin{longtable}{@{}ll@{}}
\toprule
被積分関数 & 原始関数の1つ\\
\midrule
\endhead
$x^p\ (p\ne-1)$ & $\frac{x^{p+1}}{p+1}$\\[4pt]
$\frac1x$ & $\log|x|$\\[4pt]
$\e^{bx}$ & $\frac1b\e^{bx}$\\[4pt]
$\sin bx$ & $-\frac1b\cos bx$\\[4pt]
$\cos bx$ & $\frac1b\sin bx$\\[4pt]
$\frac1{\cos^2x}$ & $\tan x$\\[4pt]
$\frac{f'(x)}{f(x)}\quad(f(x)\ne0)$ & $\log|f(x)|$\\[4pt]
$\frac1{a^2+x^2}$ & $\frac1a\arctan\frac xa$\\[4pt]
$\frac1{\sqrt{a^2-x^2}}\quad(|x|<a)$ & $\arcsin\frac xa$\\[4pt]
\bottomrule
\end{longtable}
一般の実数$p$では,まず$x>0$でべき関数を考える.逆三角関数の記号と範囲は\chapterref{章:置換の背景}で定義した.高校の基本計算では,対応する三角置換を使って定積分を求めることもできる.


